\section{Ideas on generalizing the framework}

\subsection{Conditioning operation with input nodes}
\Joris{Consider this. To incorporate possible selection bias on input nodes, there seem to be several options:
\begin{enumerate}
  \item merge the input nodes whose value range is no longer a Cartesian product, and keep their causal interpretation
  \item merge the input nodes whose value range is no longer a Cartesian product, and remove their causal interpretation
  \item add undirected edges between cliques of input nodes whose value range is no longer a Cartesian product
\end{enumerate}
Either of the two should address the factorization / CI properties.
However, it might be that undirected edges have the advantage that they express more CI statements!
Indeed, if instead of dashing endogenous ancestors of selection nodes in the conditioning SCM operation
we would merge them we may loose CI information.

Then the causal interpretation. 
\begin{enumerate}
  \item When merging nodes, the hope is that we keep the causal interpretation.
But I am not sure. So we should see if we can think of a counterexample (an s-SCM with a given causal
interpretation that gives contradictions when we exist on keeping the causal interpretation of the
merged input nodes).
  \item When adding undirected edges, I guess we would like to restrict the causal interpretation
    to joint interventions on cliques only.
\end{enumerate}
I think for a causal interpretation we have two desiderata: (i) it is a valid \emph{causal} interpretation
(so far, I believe that requires input nodes to be variation independent); (ii) it preserves the old 
causal interpretation, or indicates which part of it is lost (using e.g. dashed nodes).

One thing to check: are the $i\sigma$-separations in a CDMG preserved when putting any types of edges between input nodes?
And are the $im$-separations in a $\sigma$-iMAG preserved when putting any types of edges between input nodes?

Perhaps it's also useful to think about the purest setting with \emph{only} input variables and selection bias.
What happens to the causal interpretation there? We could limit the Cartesian product range of the input variables
through selection on a common child to any arbitrary subset of that range. 

Regarding FCI: if we are not going to test for factorization structure of the joint range of input variables, then
perhaps the conservative option is to connect them all via undirected edges. Interestingly, in the real world, the
input variables often correspond to \emph{causally independent experimental settings before filtering}. For example,
in the yeast case, we can independently knock out any number of genes (at least hypothetically). But for some 
combinations, the yeast cells won't live, and no data is produced. If we were to assume boldly that the conditioned
iMAG has independent inputs, then this communicates that we could choose any joint setting of the input nodes and
have data. This may not be true in general. Because we are assuming that the data generating mechanism is that of 
an s-SCM, we are assuming that in the absence of any selection, the input nodes are variation independent, and hence
no edges are needed. Any edge between input nodes in the presence of selection must hence be due to selection, and
hence be undirected. But to be completely honest, we could also let FCI act as if it doesn't know whether input
nodes are adjacent or not.

Aha! Note that so far we always consider a JOINT intervention on all input nodes. The ONLY two reasons to have more
than one input node are (i) so we can consider this joint intervention to be a combination of ``independent'' 
interventions; (ii) so that we can utter more refined independence statements about $P(X \mid \doit(J))$.
Note that to be able to consider the joint intervention to be a combination of ``independent'' interventions 
we need to assume that these ``independent'' interventions commute (temporally): it does not matter in which
ordering we apply them.
}

\subsection{Too good to be true?}

Here's an idea:
\begin{enumerate}
  \item Check whether for $i\sigma$-separation any additional edges between input nodes don't make any difference.
  \item Prove that SCMs also imply a variation independence Markov property.
  \item Check that Patrick's notion of CI defaults to 'independence' if it doesn't convey information.
  \item Broaden the framework by thinking of a suitable graphical representation of variation dependence between input nodes.
  \item For conditioning, work out whether it is compatible with the graphical representation or whether it suggests refining the graphical representation.
  \item (Check whether for $im$-separation any additional undirected edges between input nodes matter for the restricted separations).
\end{enumerate}

\begin{Lem}
  Let $G$ be a CDMG. Let $\tilde{G}$ denote the same ``CDMG''\footnote{Formally it's not a CDMG because we defined those without allowing for edges between input nodes, but let's still call it a CDMG for convenience.} but with additional arbitrary (directed, bidirected) edges between all input nodes of $G$.
  For $A, B, C \subseteq V \cup J$ with $A \cap J = \emptyset$ and $J \subseteq B \cup C$ we have that
    $$A \isPerp_G B \given C \iff A \isPerp_{\tilde{G}} \given C$$
\end{Lem}
\begin{proof}
  Note that the strongly connected components in $\tilde{G}$ consist of those in $G$ and in addition possible scc's within $J$.

  Suppose $A \nPerp^\sigma_G B \given C$. 
  Then there exists a $\sigma$-open path $\pi$ between $A$ and $B \cup J$ given $C$ in $G$.
  WLOG we can assume that the only node on $\pi$ that is in $J$ is its endpoint in $B \cup J$.
  The same path must be present in $\tilde{G}$.
  Any collider on $\pi$ must be in $\Anc_G(C)$.
  Hence it must be in $\Anc_{\tilde{G}}(C)$.
  Any non-collider on $\pi$ must be unblockable in $G$, or not in $C$.
  Suppose it's unblockable in $G$.
  If it's a node in $V$ then it can only point to nodes in $V$, and hence (since the scc of the node is the same in $G$ as in $\tilde{G}$) must also be unblockable in $\tilde{G}$.
  It cannot be a node in $J$ because then it must be an end node of $\pi$ and it cannot point to a node in the same scc in $G$.
  So any non-collider on $\pi$ that is unblockable in $G$ is also unblockable in $\tilde{G}$.
  But this implies that any non-collider on $\pi$ is unblockable in $\tilde{G}$, or not in $C$.
  So $\pi$ is a $\sigma$-open path between $A$ and $B \cup J$ given $C$ in $\tilde{G}$ as well.

  Suppose $A \nPerp^\sigma_{\tilde{G}} B \given C$. 
  Then there exists a $\sigma$-open path $\tilde{\pi}$ between $A$ and $B \cup J$ given $C$ in $\tilde{G}$.
  WLOG we can assume that the only node on $\tilde{\pi}$ that is in $J$ is its endpoint in $B \cup J$.
  The same path must be present in $G$.
  Any collider on $\tilde{\pi}$ must be in $V \cap \Anc_{\tilde{G}}(C)$.
  Hence it must be in $\Anc_G(C)$.
  Any non-collider on $\tilde{\pi}$ must be unblockable in $\tilde{G}$, or not in $C$.
  Suppose it's unblockable in $\tilde{G}$.
  If it's a node in $V$ then it can only point to nodes in the same scc according to $\tilde{G}$, and hence, according to $G$.
  It cannot be a node in $J$, because then it must be an endnode of $\tilde{\pi}$ and it cannot point to a node in the same scc in $\tilde{G}$.
  So any non-collider on $\tilde{\pi}$ that is unblockable in $\tilde{G}$ is also unblockable in $G$.
  But this implies that any non-collider on $\tilde{\pi}$ is unblockable in $G$, or not in $C$.
  So $\tilde{\pi}$ is a $\sigma$-open path between $A$ and $B \cup J$ given $C$ in $G$ as well.

  (TODO: clean up this proof because the two directions turned out to be identical proofs)
\end{proof}

Let's see if Patrick's notion of CI defaults to `independence' for ``non-meaningful'' cases. 
First, let's make a list of those cases.
I think we can look at the single-node CI statements that are not of restricted form.
\begin{itemize}
  \item $j_1 \idPerp j_2 \given C$: important! this is the heart of the matter!
    $$X_{j_1} \Indep X_{j_2} \given X_C \iff \exists Q(X_{j_1} \given X_C):$$
    $$ P(X_{j_1},X_{j_2},X_C \given \doit(X_J)) = Q(X_{j_1} \given X_C) \otimes P(X_{j_2},C \given \doit(X_J))$$
    Too bad! It doesn't work (take e.g.\ $C=\emptyset$).
    This might only work if we allow $Q$ to also depend on $\doit(X_J)$.
  \item $j_1 \idPerp v_1 \given C$: ignore at first (we could define it to be equivalent to $v_1 \idPerp j_1 \given C$)
  \item $v_1 \idPerp j_1 \given C$: it's not clear whether we need this in unrestricted form? (similar as in JCI paper?)
  \item $v_1 \idPerp v_2 \given C$: it's not clear whether we need this in unrestricted form? (similar as in JCI paper?)
\end{itemize}

So the only remaining option is to prove a node-dependent Markov property:
$X_{j_1} \Indep X_{j_2} \given C$ should be read as variation independence, while CI statements in restricted form should be read as Patrick's notion.
This looks promising for FCI purposes, but is not really very elegant.
(Much more elegant would be a unified notion of independence that reduces to these special cases but does not explicitly look at the node type!)

\subsubsection{New idea}

Here's yet another possibility, which is relatively straightforward to implement.
Define an sSCM by a tuple $\lt M, S, \C{S}, T, \C{T} \rt$ where the idea is that this models an SCM $M$ plus filtering mechanisms $X_S \in \C{S}$ (with $S \subseteq V$) and $X_T \in \C{T}$ (with $T \subseteq J$) (perhaps even multiple ones).
Or perhaps we can even introduce new endogenous nodes $S$ which are selection indicators.
They can be children of input nodes only, or more generally, of any subset of the nodes in $V \cup W \cup J$.

We can define its graph by using the factorization structure in $\C{T}$ to define edges between the input nodes.
This is an ideal basis for doing FCI.
The sSCM can be mapped to an SCM by performing the conditioning operation with input nodes.
Its graph can also directly be mapped to a MAG and subsequently to a PAG.
